\section{$\Bun_G$ 上的层}\srcpage{19}

为与我们由函数域 Langlands 猜想所驱动的主线保持一致，我们先在 étale 情形下定义 $\Bun_G$ 上正确的层范畴。因此，$k$ 是某个代数闭域（常见的是 $k=\overline{\F}_q$，但也可以允许 $k$ 的特征为 $0$），$e=\Ql$，而我们的层论为
\[
S\mapsto \Ind D^b_c(S,\Ql).
\]
\srcpage{20}其中 $D^b_c(S,\Ql)$ 表示可构造 étale $\Ql$-层的有界导出范畴。这里，我们把它视为稳定 $(\infty,1)$-范畴，再过渡到 Ind-对象，得到可表现 $\Ql$-线性 $(\infty,1)$-范畴。

六函子形式化的标准选择在于采用通常的张量 $\otimes$、拉回 $f^*$ 与紧支撑上同调 $f_!$-函子（及其右伴随）。在 Arinkin、Gaitsgory、Kazhdan、Raskin、Rozenblyum 与 Varshavsky（2020b）中作了不同的选择，即采用 !-张量 $\otimes^!$、例外拉回 $f^!$ 与通常上同调 $f_*$（及其右伴随）。注意
\[
D^b_c(S,\Ql)
\]
是自对偶的，即经由 Verdier 对偶等价于其反范畴；通常的 $\otimes$、$f^*$ 与 $f_!$-函子保持 $D^b_c(S,\Ql)$；而其余函子 $\otimes^!$、$f^!$ 与 $f_*$ 是 $\otimes$、$f^*$ 与 $f_!$ 在 Verdier 对偶下的共轭。因此，在紧对象层面，两种选择的差别本质上只是表面上的，仅仅在于施加一次 Verdier 对偶。然而，在过渡到 Ind-范畴并把形式化推广到叠之后，两种选择一般真正不同。幸运的是，在 $\Bun_G$ 上，这一选择仍然无关紧要：即这是光滑 Artin 叠，故 $D(\Bun_G)$ 经由光滑图谱的下降来定义。但对光滑映射 $f$，映射 $f^*$ 与 $f^!$ 至多相差一个平移后一致，所以所得的范畴与形式化的选取无关。

\begin{remark}
我必须承认，我觉得 $(-^!,-_*)$-形式化令人困惑，以至于对 Arinkin、Gaitsgory、Kazhdan、Raskin、Rozenblyum 与 Varshavsky（2020b、2020a、2021）这一系列论文，我只有表面的理解。然而，他们的选择也有清楚的动机：即与 D-模的六函子相容。

另一方面，就几何 Langlands 等价而言，这一选择并不重要。为了我个人方便，本报告将采用 $(-^*,-_!)$-形式化。
\end{remark}

\begin{remark}
与任一六函子形式化 $D$ 相联系的一个关键构造，是核的（对称幺半可表现 $(\infty,2)$-）范畴 $K_D$。定义它的方式有好几种；最简洁的一种是取 $D$ 在
\[
\Fun(\Corr(\Sch^{\mathrm{ft}}_k),\PrL_e)
\]
中的模（注意：松弛对称幺半函子恰是关于 Day 卷积的交换代数对象）。它由有限型 $k$-概形 $X$ 对应的对象 $[X]\in K_D$ 生成，且
\[
\Hom_{K_D}([X],[Y])=D(X\times_k Y),
\]
其中单位态射由 $\Delta_!\mathbf{1}$ 给出，复合由卷积给出。

在任一 2-范畴中，都有伴随态射的概念 —— 若存在态射 $\epsilon\colon\id_Y\to GF$（在 $\End(Y)$ 中）与 $\eta\colon FG\to\id_X$（在 $\End(X)$ 中），\srcpage{21}使得复合
\[
F\xrightarrow{F\epsilon}FGF\xrightarrow{\eta F}F,\qquad
G\xrightarrow{\epsilon G}GFG\xrightarrow{G\eta}G
\]
等价于恒等，则 $F\colon X\to Y$ 是 $G\colon Y\to X$ 的左伴随。伴随在存在时于同构意义下唯一。给定层 $A\in D(X)$，可将 $A$ 视为 $K_D$ 中的态射 $[*]\to[X]$，并询问左（分别右）伴随是否存在；这引出了「suave」（分别「prim」）层的概念；参见 Heyer 与 Mann（2024）。「suave」与经典意义上的泛局部无环概念密切相关，例如参见 Hansen 与 Scholze（2023）。\footnote{对常值层 $A=\mathbf{1}$ 而言，「suave」是光滑性的减弱，且与泛无环性相关，故其名字中的字母 s、u、a 甚为贴切。「prim」则是固有性的减弱。}
\end{remark}

\begin{remark}
尽管 $(-^*,-_!)$ 与 $(-^!,-_*)$ 两种形式化在 $(\infty,1)$-范畴层面大体一致，它们给出的核 $(\infty,2)$-范畴却很不相同；特别地，比较由此得到的伴随概念颇为微妙。值得注意的是，在下一个定理中，用 prim 性给出的刻画只适用于 $(-^*,-_!)$-形式化。
\end{remark}

对几何 Langlands 对应，必须限制到 $D(\Bun_G)$ 的某个子范畴。该子范畴由下述定理刻画。为陈述它，回忆：光滑概形或叠 $X$ 上层的奇异支撑是余切丛 $T^*X$ 的一个子集，衡量该层的奇异性。当 $X=\Bun_G$ 时，余切丛 $T^*X$ 成为 Higgs 丛的模空间 $\Higgs_G$，而幂零奇异支撑指的是相应的 Higgs 丛幂零，即位于 Hitchin 纤维化的零纤维中。对 étale 层，奇异支撑的概念由 Beilinson（2016）建立。

与 prim 对象的等价，其动机来自 Hansen--Mann 在 Fargues--Fontaine 曲线情形中的类似想法；它在 Kuppel 的硕士论文（2024，定理 3.3.24）中得到证明，但以 Arinkin、Gaitsgory、Kazhdan、Raskin、Rozenblyum 与 Varshavsky（2020b，猜想 14.1.8）为前提，而该猜想现已被 Gaitsgory 与 Raskin（2025a）证明。

\begin{theorem}[Arinkin、Gaitsgory、Kazhdan、Raskin、Rozenblyum 与 Varshavsky，2020b]
设 $p$ 相对于 $G$ 充分大；特别地，它对 $G$ 是非挠的素数。\footnote{这一假设是为了在 $\gc$ 中有良好的幂零锥概念。}设 $A\in D(\Bun_G)$。以下条件等价。

(1) 层 $A$ 具有幂零奇异支撑，即其奇异支撑落在 $T^*\Bun_G$ 的幂零锥中。

(2) 层 $A$ 是 Hecke-lisse 的，即对所有 $V\in\Rep(\Gd)$，Hecke 平移 $T_V(A)\in D(\Bun_G\times X)$ 沿 $X$-方向局部常值。

\srcpage{22}(3) 层 $A$ 属于
\[
\Ind D(\Bun_G)^{\prim}\subset D(\Bun_G),
\]
其中 $D(\Bun_G)^{\prim}$ 表示 prim 对象。

记
\[
D_{\Nilp}(\Bun_G)\subset D(\Bun_G)
\]
为满足这些等价条件的对象构成的全子范畴。
\end{theorem}

\begin{proof}[证明概要]
容易的方向是 (1) 与 (3) 各自推出 (2)。对 (1)，用到奇异支撑在光滑拉回与固有前推下表现良好，并借助关于幂零锥行为的一个计算来完成。参见 Nadler 与 Yun（2019，定理 5.2.1、6.1.1）。

对 (3)，先指出 prim 层自动是紧对象（Heyer 与 Mann，2024，引理 4.4.18(ii)），故该 Ind-范畴实际上是全子范畴。为验证 (3) 推出 (2)，便可归约到 prim 层。现在用到 prim 层在固有 $*$-拉回与 suave $!$-前推下稳定。这表明 $T_V(A)\in D(\Bun_G\times X)$ 必定仍是 prim 的。特别地，对任一可构造的 $B\in D(\Bun_G)$ —— 它此时是 suave 的，从而作为方向 $D(\Bun_G)\to D(*)$ 上的核具有右伴随 —— 对象 $p_!(B\otimes T_V(A))\in D(X)$ 是 prim 的。但对固有的 $X$，prim 性等价于可对偶性（Heyer 与 Mann，2024，推论 4.5.18），所以这是 $X$ 上的局部系统。若取任一闭点 $s\colon *\to\Bun_G$ 并对 $B=s_!\Ql$ 施以此论，便知 $s^*T_V(A)\in D(X)$ 是局部系统。在 $k$ 中作任意基变换后再施以同样的论证，表明 $T_V(A)$ 在 $\Bun_G$ 各点处的所有茎都是 $X$ 上的局部系统；而这正是 (2)。

由 (2) 到 (1) 的逆向蕴涵在 Arinkin、Gaitsgory、Kazhdan、Raskin、Rozenblyum 与 Varshavsky（2020b，定理 14.4.4）中建立。大致说来，要证明的是：若存在某个非幂零的奇异支撑，则通过适当选取的 Hecke 算子，会把它显现为曲线方向上的某个非平凡奇异支撑。

至于 (2) 到 (3) 的蕴涵，参见第 7 节。
\end{proof}

最后，我们说明在 Betti 与 de Rham 情形中使用的是 $\Bun_G$ 上的哪些层范畴。Betti 情形曾在 Ben-Zvi 与 Nadler（2018）中考虑。此处 $k=\C$，层论为
\[
D(S)=D(S(\C),e),
\]
即局部紧拓扑空间 $S(\C)$ 上 $e$-向量空间层的复形。同样取 $D_{\Nilp}(\Bun_G)\subset D(\Bun_G)$，即具有幂零奇异支撑的 Betti 层构成的全子范畴。同样可以证明它等价于 Hecke-lisse；但在 Betti 情形，它并不蕴涵 (ind-)prim。作为玩具例子，\srcpage{23}把 $\Bun_G$ 换成 $S^1$，并以零奇异支撑代替幂零奇异支撑：\footnote{正如我们将在第 6 节讨论的，$G=\Gm$ 的情形实质上就产生这一情形。}于是得到由 $S^1$ 上局部常值层构成的全子范畴
\[
D(*/\Z)\subset D(S^1).
\]
现在 $D(*/\Z)\cong D_{\qc}(\mathbb{G}_{m,e})$ 等价于 $\mathbb{G}_{m,e}$ 上的拟凝聚层 —— 两边都是带一个自同构的 $D(e)$ 的对象。$\mathbb{G}_{m,e}$ 的结构层给出一个紧对象，它在 $*/\Z$ 上、从而通过拉回在 $S^1$ 上，是茎为无穷维的层；这不是 (ind-)prim 的。

Betti 与 étale 情形的差异，是六函子形式化精确定义的产物：$\Ind D^b_c(-,\Ql)$ 的类比物毋宁是 $\Ind D_{\mathrm{cons}}(-,e)$，即 Zariski-可构造 Betti 层的 Ind-范畴。这一选择将与 étale 情形完全类似。

de Rham 情形更容易：此处 $k$ 是任意的特征 $0$（代数闭）域，系数 $e=k$ 与 $k$ 一致，并取 $\Bun_G$ 上 D-模的完整导出范畴 $D(\Bun_G)$。值得注意的是，并不过渡到 (ind-)holonomic D-模的某个子范畴。
